problem:  
Let \(M^n \subset N^{n+k}(c)\) be a compact, oriented, immersed submanifold in a simply connected space form of sectional curvature \(c\).  
Suppose \(M\) admits a **unit parallel normal vector field** \(V\in \Gamma(\nu(M))\), and that the corresponding **mean curvature component**
\[
H_V := \langle H, V \rangle 
\]
is constant on \(M\).  

We say such a submanifold is of **type CPMC–V** (constant parallel mean curvature in the direction \(V\)). When \(k=1\), this reduces to the ordinary constant mean curvature (CMC) condition.

It is known that totally umbilical spheres and standard products (e.g. Clifford hypersurfaces, or extrinsic products of lower-dimensional spheres and Euclidean factors) satisfy this property for some choice of \(V\). However, no classification is known in higher codimension.

**Question:**  
For a fixed codimension \(k \ge 2\), does there exist a dimension‑independent constant \(\lambda_k > 0\) such that the following rigidity property holds?

If \(M^n \subset N^{n+k}(c)\) is a compact embedded CPMC–V submanifold satisfying
\[
\|A^\circ\|^2 \le \lambda_k\, |H_V|^2  \quad \text{on } M,
\]
then \(M\) is congruent to one of the classical parallel mean curvature examples (totally umbilical, or standard spherical products inheriting parallel normals).  

Equivalently: is there a *universal extrinsic pinching threshold* forcing a constant–parallel–mean–curvature submanifold in a space form to be a standard model whenever the traceless second fundamental form stays sufficiently small compared to the \(V\)-component of mean curvature?

Why is it a "good" problem:  
This formulation corrects the ill-defined aspects of the earlier “V-minimal” concept while keeping the intended focus: **extrinsic pinching rigidity in higher codimension** with a distinguished parallel normal direction.  
The condition of constant \(H_V\) with \(\nabla^\perp V = 0\) is a standard and meaningful geometric setup—it generalizes parallel mean curvature (PMC) submanifolds and naturally extends constant mean curvature theory to higher codimensions.  

The open question about a dimension‑independent curvature pinching constant \(\lambda_k\) touches deep territory: it asks whether quantitative control on the traceless second fundamental form relative to one mean curvature component forces rigidity, a phenomenon classically known only in codimension one and under stronger curvature assumptions.  

An affirmative answer would generalize the celebrated Simons‑type “gap” results to the realm of parallel‑mean‑curvature submanifolds with a preferred normal direction, revealing new interactions between **extrinsic curvature**, **codimension geometry**, and **analytic curvature identities**. Conversely, any counterexample would expose fundamentally new stable configurations in higher codimension, expanding our understanding of submanifold geometry in space forms.  

Thus the problem remains concrete, well‑posed, and deeply linked to major themes in modern differential geometry, while staying open and potentially groundbreaking.